% chapters/ch10_control_pid.tex — Chapter 10: Control Systems & PID
% Scope (PLAN.md): open vs closed loop, P/I/D intuition, tuning incl.
% Ziegler-Nichols, overshoot/rise/settling/steady-state error, sensor feedback
% & noise filtering, discrete PID in C, servo position-loop walkthrough.
% Budget 30 pp, mix 20 Q&A / 6 code / 4 concept.
% All snippets verified gcc -Wall -Wextra -std=c11 -lm (zero warnings).

\chapter{Control Systems \& PID}
\label{ch:control-pid}

\begin{partnote}
\textbf{Scope} --- open vs closed loop $\rightarrow$ the P, I, and D terms
(what each fixes and breaks) $\rightarrow$ tuning, including Ziegler--Nichols
$\rightarrow$ overshoot / rise time / settling time / steady-state error
$\rightarrow$ sensor feedback \& noise filtering $\rightarrow$ a discrete PID in
C $\rightarrow$ a servo position-loop walkthrough. Control theory underpins
closed-loop actuator and servo work; this chapter is answerable at the depth a
hardware/embedded interviewer expects --- intuition plus a working discrete
controller --- without a controls-PhD detour.
\end{partnote}

A PID controller is the most-deployed control algorithm in the world, and the
most-asked control question in an embedded interview. The goal here is not to
derive transfer functions but to \emph{reason} about a loop: why it oscillates,
why it never quite reaches target, what each gain does, and how to write and
tune one in C on a microcontroller.

%=====================================================================
\section{Primer}
%=====================================================================

\subsection{Open loop vs closed loop}

\textbf{Open loop} commands an actuator with no measurement of the result ---
``run the motor at 60\% duty'' --- simple, but it can't correct for disturbances
(load, friction, supply sag) and has no idea whether it hit the target.
\textbf{Closed loop} measures the output, computes the \textbf{error}
(setpoint $-$ measurement), and feeds it back to drive the actuator so the error
shrinks. Feedback buys disturbance rejection and accuracy at the cost of
complexity and the risk of \emph{instability} (a badly-tuned loop oscillates or
diverges). Every servo position loop, temperature controller, and motor-speed
regulator is closed loop.

\subsection{The three terms}

PID sums three responses to the error $e$:
\[
u(t) = K_p\,e + K_i\!\int e\,dt + K_d\,\frac{de}{dt}.
\]
\textbf{Proportional} ($K_p$) reacts to the \emph{present} error --- bigger
error, bigger push. Alone it leaves a \emph{steady-state error} (it needs some
error to produce output) and, pushed high, causes overshoot and oscillation.
\textbf{Integral} ($K_i$) accumulates \emph{past} error, driving the
steady-state error to zero --- but it adds lag and can \emph{wind up}
(over-accumulate) causing large overshoot. \textbf{Derivative} ($K_d$)
anticipates \emph{future} error from its rate of change, adding damping that
reduces overshoot and settling time --- but it amplifies sensor noise. The art
is balancing them.

\subsection{Response metrics}

A loop's quality is judged by its \textbf{step response} (how it reacts when the
setpoint jumps): \textbf{rise time} (how fast it first reaches the target),
\textbf{overshoot} (how far it exceeds target, as a percentage),
\textbf{settling time} (how long until it stays within a band, e.g.\ $\pm$2\%,
of target), and \textbf{steady-state error} (the residual offset once settled).
These trade off: more aggressive (faster rise) usually means more overshoot, so
tuning is choosing the right compromise for the application --- a camera gimbal
wants no overshoot, a fast positioner may tolerate a little.

\subsection{Discretisation, noise, and windup}

On an MCU the loop runs at a fixed \textbf{sample period} $dt$ (the timer-driven
control frequency from Chapter 4), so the integral becomes a running sum
($\times dt$) and the derivative a difference ($/dt$). Three practical
must-dos: (1) \textbf{filter the measurement} (a low-pass / moving average)
before the derivative, or noise is amplified into the output; (2) \textbf{take
the derivative on the measurement, not the error}, to avoid a ``derivative
kick'' spike when the setpoint steps; (3) \textbf{anti-windup} --- stop the
integral accumulating while the actuator is saturated, or it overshoots badly on
recovery. A correct embedded PID is these three fixes plus the three terms.

%=====================================================================
\section{Q\&A Bank}
%=====================================================================

\tierbanner{Tier 1 --- Screening (phone / HR-technical)}%
{Two to four sentences.}

\question{Q1.1}{Open loop vs closed loop?}
Open loop commands the actuator without measuring the result, so it can't
correct for disturbances; closed loop measures the output, computes the error
versus the setpoint, and feeds it back to drive the error to zero. Closed loop
gives accuracy and disturbance rejection but can be unstable if poorly tuned.

\question{Q1.2}{What does the P term do?}
Proportional control outputs in proportion to the present error --- larger
error, larger correction. It's the main driving term, but alone it leaves a
residual steady-state error and, if too high, causes overshoot and oscillation.

\question{Q1.3}{What does the I term do?}
Integral accumulates past error over time, so it keeps pushing until the error
is exactly zero --- eliminating steady-state error. Its downside is added lag and
integral windup, which cause overshoot.

\question{Q1.4}{What does the D term do?}
Derivative responds to the rate of change of error, anticipating where it's
heading and adding damping --- reducing overshoot and settling time. It amplifies
measurement noise, so it usually needs a filtered input.

\question{Q1.5}{What is overshoot?}
How far the response exceeds the target before settling, usually expressed as a
percentage of the setpoint. Too much overshoot can damage mechanisms or
destabilise the system; some applications require zero overshoot.

\question{Q1.6}{What is steady-state error?}
The residual difference between the setpoint and the actual settled output. A
pure-P controller leaves a steady-state error; adding integral action removes
it.

\question{Q1.7}{What is settling time?}
The time for the response to enter and stay within a defined band (commonly
$\pm$2\% or $\pm$5\%) of the final value. It measures how quickly the loop
``finishes'' reacting to a change.

\question{Q1.8}{Why filter a sensor before the derivative term?}
Because the derivative amplifies high-frequency content, so raw sensor noise
becomes large, jittery output that can buzz the actuator. A low-pass filter (or
moving average) smooths the measurement so the derivative reflects real motion,
not noise.

\question{Q1.9}{What is integral windup?}
When the actuator saturates (can't move further) but the integral keeps
accumulating error, building a large value that causes a big overshoot when the
error finally reverses. Anti-windup limits or freezes the integral while
saturated.

\question{Q1.10}{What sets the control loop's sample rate?}
A fixed periodic timer interrupt (Chapter 4) runs the PID at a constant $dt$;
the rate must be fast enough relative to the plant's dynamics (a common rule of
thumb is 10--20$\times$ the loop bandwidth) and have low jitter, since the
integral and derivative math assume a constant period.

\tierbanner{Tier 2 --- Core technical (panel round)}%
{A short paragraph.}

\question{Q2.1}{Walk through tuning a PID from scratch.}
Start with $K_i = K_d = 0$ and raise $K_p$ until the response is fast but begins
to oscillate; back it off slightly. Add $K_i$ to remove the steady-state error,
increasing it until the offset is eliminated without too much overshoot/
sluggishness. Add $K_d$ to damp the overshoot and shorten settling, increasing
until noise starts to show in the output, then back off. Iterate. This manual
method (or a structured one like Ziegler--Nichols) trades the classic
fast-vs-stable compromise; you tune to the application's priority (no overshoot
vs fast response).

\question{Q2.2}{Explain Ziegler--Nichols tuning.}
A structured method: with only P active, raise $K_p$ until the loop oscillates
with a steady amplitude --- that gain is the \textbf{ultimate gain} $K_u$ and the
oscillation period is $T_u$. Then set the PID gains from standard ratios: for a
classic PID, $K_p = 0.6K_u$, $K_i = 1.2K_u/T_u$, $K_d = 0.075K_uT_u$ (equivalently
$T_i = T_u/2$, $T_d = T_u/8$). It gives a quick, repeatable starting point ---
typically a bit aggressive (~25\% overshoot) --- which you then refine. Its
limitation is that driving a real system into sustained oscillation may be
unsafe, so it's not always usable on flight hardware.

\question{Q2.3}{Why does a proportional-only controller have steady-state
error?}
Because its output is $K_p \times$ error, so to hold any non-zero output
(needed to counter gravity, friction, or a load) it \emph{requires} a non-zero
error --- if the error were zero, the output would be zero and the system would
sag. The error settles at exactly the value whose $K_p$ product holds the load.
Integral action fixes this: it keeps integrating that residual error until its
own term supplies the needed output, allowing the error to reach zero.

\question{Q2.4}{What is derivative kick and how do you avoid it?}
When the setpoint changes in a step, the error jumps instantly, so its
derivative is a huge spike --- the D term produces a violent output transient
(``kick'') that can slam the actuator. The fix is to compute the derivative on
the \emph{measurement} instead of the error: since the measured value can't jump
instantaneously, there's no spike on a setpoint change, while the damping
behaviour during motion is unchanged (you negate it, because $d(\text{error})/dt
= -d(\text{measurement})/dt$ when the setpoint is constant).

\question{Q2.5}{Describe a servo position control loop.}
The setpoint is the desired angle/position; a sensor (encoder, potentiometer,
resolver) measures the actual position; the PID computes a command from the
error; that command becomes a PWM duty (Chapter 4) driving the motor/actuator,
which moves the load, changing the measured position --- closing the loop. Often
it's a cascaded loop: an outer position loop sets the target for an inner
velocity (and sometimes current) loop, each tuned for its own bandwidth. This is
the closed-loop actuator structure behind a CAN servo (Chapter 7).

\question{Q2.6}{How do you implement anti-windup?}
Two common methods. \textbf{Clamping/conditional integration}: only accumulate
into the integral when the output is \emph{not} saturated (or when integrating
would move the output away from the saturated rail), so the integral can't build
up while the actuator is maxed out. \textbf{Back-calculation}: feed the
difference between the saturated and unsaturated output back to bleed the
integral down. Either way the integral reflects achievable output, so the loop
recovers cleanly from saturation without a huge overshoot --- essential whenever
the actuator has limited range.

\question{Q2.7}{How do you choose the control loop frequency?}
Fast enough that the plant looks ``continuous'' to the controller --- a common
guideline is 10--20$\times$ the desired closed-loop bandwidth (or the fastest
plant time constant), and fast relative to disturbance frequencies. But faster
costs CPU and amplifies derivative noise, and the sensor's own update rate and
the actuator's response set a ceiling. You also need \emph{low jitter}
(Chapter 4), because the discrete integral/derivative assume a constant $dt$;
a 1\,kHz timer-driven loop is typical for a servo.

\question{Q2.8}{Why must the sample period be constant (low jitter)?}
Because the discrete PID multiplies the integral by $dt$ and divides the
derivative by $dt$ assuming a fixed period; if $dt$ varies frame to frame, the
effective $K_i$ and $K_d$ wander, degrading stability and adding noise to the
output. That's why the loop runs from a dedicated periodic timer interrupt with
high priority and short execution, not from a software-delay loop --- tying
directly to the real-time determinism of Chapters 4 and 9.

\tierbanner{Tier 3 --- Trap \& depth (principal-engineer level)}%
{A paragraph plus the trap.}

\question{Q3.1}{Your loop oscillates. Walk through diagnosing whether it's too
much P, too much I, too little D, or something physical.}
Read the \emph{character} of the oscillation. Fast oscillation that grows with
$K_p$ points to excess proportional gain (reduce $K_p$ or add $K_d$ damping). A
slow oscillation/hunting that appears when you added integral suggests too much
$K_i$ (or windup) --- the integral overshoots, reverses, overshoots. Oscillation
that's actually high-frequency jitter on the output usually means $K_d$
amplifying sensor noise (filter the measurement, reduce $K_d$). But also rule out
\emph{physical} causes: mechanical backlash/compliance, actuator saturation
cycling, a sample rate too slow for the plant, or transport delay (a pure time
delay makes \emph{any} high-gain loop unstable). \trap the trap is reflexively
turning knobs; the disciplined approach identifies the oscillation frequency and
amplitude (scope the measurement), relates it to which term or physical effect
produces that signature, and checks $dt$/saturation/delay before re-tuning ---
the same instrument-led method as the protocol chapters.

\question{Q3.2}{Integral action removes steady-state error --- so why not just
use a huge $K_i$?}
Because integral adds \emph{phase lag}: it responds to accumulated past error,
so it's always ``behind,'' and too much of it makes the loop sluggish and
oscillatory, and it winds up badly against saturation. A large $K_i$ removes the
offset \emph{faster} but at the cost of large overshoot and reduced stability
margin --- it can even destabilise the loop. The right amount is the minimum that
eliminates the offset within the required settling time, paired with anti-windup.
\trap the trap is treating ``$I$ kills steady-state error'' as ``more $I$ is
better.'' Each term has a cost: $P$ $\rightarrow$ overshoot, $I$ $\rightarrow$
lag/windup, $D$ $\rightarrow$ noise. Senior answers name the \emph{cost} of each
term, not just its benefit.

\question{Q3.3}{When is a PID the wrong tool, and what would you reach for
instead?}
PID assumes a roughly linear, single-input/single-output plant with modest
delay. It struggles with: significant \textbf{transport delay} (use a Smith
predictor or model-based control), strongly \textbf{nonlinear} plants (gain
scheduling --- different PID gains per operating point --- or feedback
linearisation), \textbf{coupled multi-variable} systems (state-space / LQR /
MPC), and cases where you can \emph{model} the disturbance (feed-forward added to
the PID dramatically improves tracking). For a flight control law specifically,
the design is typically state-space/gain-scheduled, not a lone PID. \trap the
trap is forcing PID onto everything; the depth is knowing its assumptions and
naming the right alternative --- and that adding \textbf{feed-forward} (command
the known steady output, let PID trim the rest) often beats tuning PID harder.
Saying ``I verify the control law; the law design was the systems team's'' is
also the honest boundary the resume implies.

\question{Q3.4}{Two channels run the same PID on the same sensor but produce
slightly different outputs each frame. Is that a fault?}
Not necessarily --- it's expected (Chapter 8). Each channel's ADC quantises
differently, sensor noise differs, floating-point rounding and timing jitter
differ, and the integral state can diverge slightly over time even from
identical inputs. So a redundant/voting system must compare channel outputs with
a \textbf{tolerance band} and \textbf{persistence}, not exact equality, before
declaring a channel faulty --- and may periodically \emph{re-synchronise} the
integral states across channels via the CCDL to stop slow divergence. \trap the
trap is expecting bit-identical outputs from redundant PID channels and
nuisance-tripping on normal divergence; the real design bounds acceptable
divergence and resynchronises integrator state --- exactly the CCDL
integrity/voting verification on the resume, now connected to the control law it
protects.

%=====================================================================
\section{Coding Gym}
%=====================================================================

\begin{partnote}
Six host-compilable problems (verified under
\code{gcc -Wall -Wextra -std=c11 -lm}) building a production-quality discrete
PID step by step: the basic controller, anti-windup, derivative-on-measurement,
the sensor filter, Ziegler--Nichols gains, and step-response metrics.
\end{partnote}

%---------------------------------------------------------------
\begin{gymproblem}{Discrete PID controller}

\textbf{Problem.} Implement one step of a positional discrete PID given
setpoint, measurement, and sample period.

\textbf{Hints.}
\begin{itemize}
\item Integral $+= e\cdot dt$; derivative $=(e-e_{\text{prev}})/dt$; output is
  the weighted sum.
\end{itemize}

\attempt

\textbf{Solution.}
\begin{cgym}
#include <stdio.h>

typedef struct { double kp, ki, kd, integ, prev_err; } pid_t;

static double pid_step(pid_t *p, double setpoint, double meas, double dt) {
    double err   = setpoint - meas;
    p->integ    += err * dt;                  /* accumulate past error  */
    double deriv = (err - p->prev_err) / dt;  /* rate of change         */
    p->prev_err  = err;
    return p->kp * err + p->ki * p->integ + p->kd * deriv;
}

int main(void) {
    pid_t p = {2.0, 0.5, 0.1, 0, 0};
    printf("u=%.3f\n", pid_step(&p, 1.0, 0.0, 0.01));   /* 12.005 */
    return 0;
}
\end{cgym}

This is the textbook positional form: each call computes the error, updates the
integral and derivative using the fixed $dt$, and returns the actuator command.
It runs from the periodic control-loop timer (Chapter 4). The next three problems
turn this naive version into one that survives real hardware --- saturation,
setpoint steps, and sensor noise.

\followups
\begin{enumerate}
\item \emph{``Velocity (incremental) form?''} $\rightarrow$ Compute $\Delta u$
  from successive errors; useful when the actuator integrates (e.g.\ a stepper).
\item \emph{``What are the units of the output?''} $\rightarrow$ Whatever the
  actuator takes --- scale/clamp to PWM duty or current command.
\item \emph{``Why positional vs velocity form?''} $\rightarrow$ Positional is
  intuitive; velocity form has natural bumpless transfer and simpler windup
  handling.
\end{enumerate}
\end{gymproblem}

%---------------------------------------------------------------
\begin{gymproblem}{PID with anti-windup and output clamping}

\textbf{Problem.} Add output saturation limits and stop the integral winding up
while saturated.

\textbf{Hints.}
\begin{itemize}
\item Only integrate when the (unsaturated) output is within limits; clamp the
  final output.
\end{itemize}

\attempt

\textbf{Solution.}
\begin{cgym}
#include <stdio.h>

typedef struct { double kp, ki, kd, integ, prev_err, omin, omax; } pidaw_t;

static double pid_aw_step(pidaw_t *p, double sp, double meas, double dt) {
    double err   = sp - meas;
    double deriv = (err - p->prev_err) / dt;
    p->prev_err  = err;
    double out_unsat = p->kp * err + p->ki * (p->integ + err * dt) + p->kd * deriv;
    if (out_unsat <= p->omax && out_unsat >= p->omin)   /* conditional integration */
        p->integ += err * dt;
    double out = p->kp * err + p->ki * p->integ + p->kd * deriv;
    if (out > p->omax) out = p->omax;                   /* clamp */
    else if (out < p->omin) out = p->omin;
    return out;
}

int main(void) {
    pidaw_t p = {2, 0.5, 0.1, 0, 0, -1.0, 1.0};
    double o = 0;
    for (int i = 0; i < 5; i++) o = pid_aw_step(&p, 10.0, 0.0, 0.01);
    printf("out=%.3f (clamped to 1.0)\n", o);   /* 1.000 */
    return 0;
}
\end{cgym}

With a large error (setpoint 10, output limited to $\pm$1) the actuator
saturates; \emph{conditional integration} freezes the integral while saturated
so it can't wind up, and the output is clamped to the physical limit. Without
this, the integral would balloon and cause a huge overshoot when the
measurement finally catches up --- the classic windup failure (Q1.9, Q2.6).
Every real actuator has limits, so anti-windup is mandatory, not optional.

\followups
\begin{enumerate}
\item \emph{``Back-calculation anti-windup instead?''} $\rightarrow$ Bleed the
  integral by $(out_{\text{sat}} - out_{\text{unsat}})\cdot K_{aw}$.
\item \emph{``Why integrate the \emph{unsaturated} check?''} $\rightarrow$ To
  decide whether adding more integral would push further into saturation.
\item \emph{``Asymmetric limits?''} $\rightarrow$ \code{omin}/\code{omax} can
  differ (e.g.\ a motor that brakes harder than it drives).
\end{enumerate}
\end{gymproblem}

%---------------------------------------------------------------
\begin{gymproblem}{Derivative on measurement (no kick)}

\textbf{Problem.} Compute the derivative term from the measurement instead of
the error to avoid derivative kick on setpoint steps.

\textbf{Hints.}
\begin{itemize}
\item $d(\text{error})/dt = -d(\text{meas})/dt$ for a constant setpoint; track
  \code{prev\_meas}.
\end{itemize}

\attempt

\textbf{Solution.}
\begin{cgym}
#include <stdio.h>

typedef struct { double kp, ki, kd, integ, prev_meas; } piddm_t;

static double pid_dm_step(piddm_t *p, double sp, double meas, double dt) {
    double err   = sp - meas;
    p->integ    += err * dt;
    double deriv = -(meas - p->prev_meas) / dt;   /* derivative of measurement */
    p->prev_meas = meas;
    return p->kp * err + p->ki * p->integ + p->kd * deriv;
}

int main(void) {
    piddm_t p = {2, 0.5, 0.1, 0, 0};
    printf("u=%.3f\n", pid_dm_step(&p, 1.0, 0.0, 0.01));   /* 2.005 */
    return 0;
}
\end{cgym}

When the operator steps the setpoint, error jumps but the \emph{measurement}
doesn't, so deriving from the measurement gives no spike --- no violent actuator
``kick'' (Q2.4). During actual motion the measurement changes smoothly and the
damping is identical, because with a constant setpoint the two derivatives differ
only in sign (hence the negation). Compare the output here (2.005) with the
basic version's (12.005) on the same setpoint step: the kick is gone. This is a
standard feature of industrial PID blocks.

\followups
\begin{enumerate}
\item \emph{``Derivative on error vs measurement during tracking --- same?''}
  $\rightarrow$ Identical for constant setpoint; only setpoint changes differ.
\item \emph{``Combine with a filter on \code{meas}.''} $\rightarrow$ Yes ---
  filter then differentiate (next problem) to also kill noise.
\item \emph{``Setpoint weighting?''} $\rightarrow$ A generalisation that scales
  how much P/D see the setpoint vs measurement.
\end{enumerate}
\end{gymproblem}

%---------------------------------------------------------------
\begin{gymproblem}{Low-pass filter the sensor (EMA)}

\textbf{Problem.} Implement an exponential moving-average filter to smooth a
noisy measurement before the derivative.

\textbf{Hints.}
\begin{itemize}
\item \code{y += alpha*(x - y)}; smaller \code{alpha} = more smoothing.
\end{itemize}

\attempt

\textbf{Solution.}
\begin{cgym}
#include <stdio.h>

static double ema(double prev, double sample, double alpha) {
    return prev + alpha * (sample - prev);   /* 0<alpha<=1; smaller = smoother */
}

int main(void) {
    double y = 0;
    double in[] = {0, 1, 1, 1, 1};
    for (int i = 0; i < 5; i++) y = ema(y, in[i], 0.5);
    printf("filtered=%.4f\n", y);   /* 0.9375 -> converging to 1 */
    return 0;
}
\end{cgym}

The EMA is the cheapest useful low-pass filter --- one multiply-add, one state
variable, no buffer (unlike a moving average). \code{alpha} sets the cutoff: near
1 follows the input closely (little filtering), near 0 heavily smooths (more lag).
You filter the measurement \emph{before} the derivative term so the D action
responds to real motion, not sensor noise (Q1.8) --- the trade is that filtering
adds lag, which can reduce stability, so you pick the lightest filter that tames
the noise.

\followups
\begin{enumerate}
\item \emph{``Relate \code{alpha} to cutoff frequency.''} $\rightarrow$
  \code{alpha = dt/(RC+dt)}; choose RC for the desired corner.
\item \emph{``Why not a long moving average?''} $\rightarrow$ More lag and RAM;
  EMA gives similar smoothing with $O(1)$ state.
\item \emph{``Filter lag hurts the loop --- mitigation?''} $\rightarrow$ Keep
  the filter light, or filter only the D path, not the whole measurement.
\end{enumerate}
\end{gymproblem}

%---------------------------------------------------------------
\begin{gymproblem}{Ziegler--Nichols gains}

\textbf{Problem.} Compute classic PID gains from the ultimate gain $K_u$ and
oscillation period $T_u$.

\textbf{Hints.}
\begin{itemize}
\item $K_p=0.6K_u$, $K_i=1.2K_u/T_u$, $K_d=0.075K_uT_u$.
\end{itemize}

\attempt

\textbf{Solution.}
\begin{cgym}
#include <stdio.h>

static void zn_pid(double Ku, double Tu, double *kp, double *ki, double *kd) {
    *kp = 0.6 * Ku;
    *ki = 1.2 * Ku / Tu;
    *kd = 0.075 * Ku * Tu;
}

int main(void) {
    double kp, ki, kd;
    zn_pid(4.0, 2.0, &kp, &ki, &kd);
    printf("kp=%.3f ki=%.3f kd=%.3f\n", kp, ki, kd);   /* 2.400 2.400 0.600 */
    return 0;
}
\end{cgym}

You find $K_u$ by raising P-only gain until the loop oscillates with constant
amplitude, and $T_u$ is that oscillation's period; the ratios then give a
working PID (Q2.2). It's a fast, repeatable \emph{starting point} --- typically a
touch aggressive (~25\% overshoot) --- that you then refine by hand. The caveat
for flight hardware: deliberately driving the system into sustained oscillation
to measure $K_u$ may be unsafe, so model-based tuning or careful step-response
methods are often used instead.

\followups
\begin{enumerate}
\item \emph{``PI-only ZN ratios?''} $\rightarrow$ $K_p=0.45K_u$,
  $K_i=0.54K_u/T_u$ --- when D-noise is a problem.
\item \emph{``Why is ZN often too aggressive?''} $\rightarrow$ It targets
  quarter-amplitude decay; many apps want less overshoot --- detune from here.
\item \emph{``Alternative that doesn't need oscillation?''} $\rightarrow$
  Step-response (Cohen--Coon) or model-based tuning.
\end{enumerate}
\end{gymproblem}

%---------------------------------------------------------------
\begin{gymproblem}{Step-response metrics}

\textbf{Problem.} From a sampled step response, compute percent overshoot and
2\% settling time.

\textbf{Hints.}
\begin{itemize}
\item Overshoot = (peak $-$ setpoint)/setpoint; settling = last time it leaves
  the $\pm$2\% band.
\end{itemize}

\attempt

\textbf{Solution.}
\begin{cgym}
#include <stdio.h>
#include <math.h>

static void step_metrics(const double *y, int n, double dt, double setpoint,
                         double *overshoot_pct, double *settling_time) {
    double peak = 0;
    for (int i = 0; i < n; i++) if (y[i] > peak) peak = y[i];
    *overshoot_pct = (peak - setpoint) / setpoint * 100.0;
    double band = 0.02 * setpoint;
    int settle_idx = n - 1;
    for (int i = n - 1; i >= 0; i--)
        if (fabs(y[i] - setpoint) > band) { settle_idx = i + 1; break; }
    *settling_time = settle_idx * dt;
}

int main(void) {
    double y[] = {0, 0.5, 1.2, 1.1, 1.02, 1.0, 1.0, 1.0};
    double os, ts;
    step_metrics(y, 8, 0.1, 1.0, &os, &ts);
    printf("overshoot=%.1f%% settling=%.1fs\n", os, ts);   /* 20.0% 0.5s */
    return 0;
}
\end{cgym}

This response peaks at 1.2 (20\% overshoot) and last leaves the $\pm$2\% band at
$t=0.5$\,s. These are the numbers you quote to characterise a tuned loop, and the
ones you watch on the scope while tuning --- pushing for faster settling usually
increases overshoot, so you tune to the application's tolerance. Computing them
on-target from a logged step lets you verify the controller meets its spec
objectively rather than by eye.

\followups
\begin{enumerate}
\item \emph{``Add rise time (10--90\%).''} $\rightarrow$ Find the times crossing
  0.1 and 0.9 of setpoint; difference is rise time.
\item \emph{``Steady-state error from this array?''} $\rightarrow$ setpoint
  minus the final settled value.
\item \emph{``Why measure rather than eyeball?''} $\rightarrow$ Objective,
  repeatable pass/fail against the requirement --- the V\&V mindset (Chapter 12).
\end{enumerate}
\end{gymproblem}

%=====================================================================
\section{Link to Her Resume}
%=====================================================================

\begin{resumelink}
\textbf{Closed-loop fundamentals:} ``My academic project was a closed-loop
PLC automation system with sensor feedback, and control systems was core
coursework. I can reason about a loop --- why it overshoots or never reaches
target, what each PID term costs, and how to write a discrete PID with
anti-windup and a filtered derivative on a microcontroller at a fixed,
low-jitter sample rate.''

\medskip
\textbf{Servo / actuator context:} ``A CAN servo is a closed-loop position
controller --- setpoint in, encoder feedback, PID to a PWM command, often
cascaded position/velocity loops. I understand that structure and how the
control frequency and jitter (from the real-time timer) affect it.''

\medskip
\textbf{Honest boundary:} ``I verify control-law \emph{outputs} and the loop's
behaviour; designing the flight control law itself is the systems team's work,
and that's typically state-space/gain-scheduled rather than a lone PID --- but I
can explain how I'd verify its response (overshoot, settling, steady-state
error) objectively.''

\medskip
\small Maps to the resume's control-systems coursework, the closed-loop
academic project, and DFCC/servo context --- with an honest scope boundary on
control-law \emph{design} vs \emph{verification}. Deep project scripting is
Chapter 15, from \code{inputs/INTAKE.md} only.
\end{resumelink}

%=====================================================================
\section{Self-Test Scorecard}
%=====================================================================

\begin{scorecard}
\begin{enumerate}
\item Open vs closed loop --- what does feedback buy and risk?
\item What does each of P, I, D respond to (present/past/future error)?
\item Why does P-only control leave a steady-state error?
\item What removes steady-state error, and what's its cost?
\item What is derivative kick and the fix?
\item What is integral windup and how do you prevent it?
\item State the Ziegler--Nichols classic-PID gain ratios from $K_u$, $T_u$.
\item Why must the control sample period be constant?
\item Name two response metrics and the trade-off between them.
\item When is a PID the wrong tool, and what replaces it?
\end{enumerate}
\end{scorecard}

\begin{answerkey}
\begin{enumerate}
\item Closed loop measures output and corrects error (accuracy + disturbance
  rejection); the risk is instability if poorly tuned. Open loop can't correct
  at all.
\item P = present error, I = accumulated past error, D = rate of change
  (anticipates future).
\item Its output is $K_p\times$error, so holding a non-zero output requires a
  non-zero error.
\item Integral action; its cost is added lag and windup (overshoot).
\item A spike from the error's derivative on a setpoint step; compute the
  derivative on the measurement instead.
\item The integral accumulating while the actuator is saturated; freeze/limit
  it (conditional integration or back-calculation).
\item $K_p=0.6K_u$, $K_i=1.2K_u/T_u$, $K_d=0.075K_uT_u$.
\item The discrete integral ($\times dt$) and derivative ($/dt$) assume a fixed
  period; jitter makes effective $K_i$/$K_d$ wander.
\item E.g.\ overshoot vs settling/rise time --- faster response generally costs
  more overshoot.
\item For large delay, strong nonlinearity, or coupled MIMO plants; use
  Smith-predictor/model-based, gain scheduling, or state-space/LQR/MPC (and add
  feed-forward).
\end{enumerate}
\end{answerkey}

\begin{partnote}
\emph{End of Chapter 10.} Next: \textbf{Hardware Debugging \& Lab Skills} ---
the oscilloscope, logic analyser, JTAG/SWD, ST-Link vs J-Link, and the
board-bring-up sequence --- the instrument-led methods this book keeps invoking,
made explicit.
\end{partnote}
